\clearpage\fancyhead[L]{\small GAUSS AN BODE.}\hypertarget{native-362-lower}{}\label{native-362-lower}
\begin{center}\textbf{GAUSS AN BODE.}\par\vspace{4pt}\hrule\vspace{4pt}\fontsize{10}{12}\selectfont Astronomisches Jahrbuch für 1815. Seite 245\;..\;249. Berlin 1812.\par\vspace{4pt}\hrule\vspace{4pt}\end{center}

\begin{flushright}Göttingen, 22. August 1812.\end{flushright}
\par\indent Wenn ich mich nicht irre, habe ich Ihnen bereits von der ersten Wiederauffindung der Pallas berichtet. Diesmal sind hier nur wenige Beobachtungen gemacht worden. Um die Zeit der Opposition war der Stand des
\SourceLeaf{363}Planeten für unseren Standort zu hoch, um ihn mit dem gewöhnlichen Instrument am Kreismikrometer zu beobachten. Ich musste deshalb ein anderes, schlechteres Instrument verwenden. Auch damit ließen schlechtes Wetter und Mondschein nur eine einzige Beobachtung zu. Hier folgen alle meine Beobachtungen:

\begin{center}\fontsize{10}{12}\selectfont\setlength{\tabcolsep}{4pt}\renewcommand{\arraystretch}{1.15}\begin{tabular}{l|c|c|c}
& Mittlere Zeit & Scheinb. Rektasz. & Nördl. Dekl.\\\hline
April 9 & $11^{\mathrm u}\;32^{\mathrm m}\;27^{\mathrm s}$ & $269^\circ\;5'\;58''4$ & $16^\circ\;38'$\\
Mai 2 & $10\;36\;6$ & $268\;42\;5.3$ & $21\;4\;31''7$\\
— 4 & $10\;44\;44$ & $268\;31\;38.0$ & $21\;25\;4.3$\\
Juni 7 & $11\;13\;51$ & $262\;50\;51.1$ & $25\;2\;56.6$\\
\end{tabular}\end{center}

\par\indent Außer zwei früheren Beobachtungen von Herrn Dr. \textsc{Olbers} wurden mir von auswärts nur vier Meridianbeobachtungen von Herrn Dr. \textsc{Burckhardt} an der Sternwarte der Militärschule in Paris mitgeteilt. Ich habe sie in Nr. 127 unserer Gelehrten Anzeigen [S. 356 dieses Bandes] abdrucken lassen. Aus diesen Daten hat Herr \textsc{Nicolai}, dessen Geschick im Rechnen Sie bereits kennen, die Opposition berechnet:

\Needspace{10\baselineskip}1812 Juni 10. $3^{\mathrm u}\;2^{\mathrm m}\;48^{\mathrm s}$ Mittlere Zeit in Göttingen

\begin{center}\fontsize{10}{12}\selectfont\setlength{\tabcolsep}{2pt}\renewcommand{\arraystretch}{1.15}\begin{tabular}{lcr}
Wahre Länge & $.\;.\;.\;.\;.\;.\;.\;.\;.\;.\;.$ & $259^\circ\;27'\;51''7$\\
Wahre geoz. Breite & $.\;.\;.\;.\;.\;.\;.\;.\;.$ & $48\;16\;26.1\;\mathrm{N}.$\\
\end{tabular}\end{center}

\Needspace{13\baselineskip}\par\indent Nachdem an den letzten Elementen nur die folgenden kleinen Korrekturen vorgenommen worden waren,

\begin{center}\fontsize{10}{12}\selectfont\setlength{\tabcolsep}{4pt}\renewcommand{\arraystretch}{1.15}\begin{tabular}{rl}
$-13''29$ & bei der Epoche der Länge 1810\\
$-0.01200$ & bei der täglichen Bewegung\\
$-17.57$ & bei der Länge des Perihels\\
$-4.09$ & beim Exzentrizitätswinkel\\
$+0''80$ & bei der Länge des Knotens\\
$-1.21$ & bei der Neigung der Bahn\\
$+0.0000045$ & beim Log. der großen Halbachse;\\
\end{tabular}\end{center}

\Needspace{17\baselineskip}\noindent ergab sich folgende größtmögliche Übereinstimmung.

\begin{center}\fontsize{10}{12}\selectfont\setlength{\tabcolsep}{6pt}\renewcommand{\arraystretch}{1.15}\begin{tabular}{c|c|c}
& \multicolumn{2}{c}{Differenz}\\
Oppos. & Mittl. Länge & helioz. Breite\\\hline
$1803$ & $-5''06$ & $-2''78$\\
$1804$ & $-12.95$ & $-9.66$\\
$1805$ & $-15.20$ & $-0.67$\\
$1807$ & $+8.28$ & $-5.77$\\
$1808$ & $+16.04$ & $+11.67$\\
$1809$ & $+13.00$ & $-1.48$\\
$1811$ & $+11.73$ & $-9.53$\\
$1812$ & $-25.92$ & $+6.15$\\
\end{tabular}\end{center}

\par\indent Ich schreibe diese Unterschiede hauptsächlich den Störungen durch Saturn und Mars zu, da bisher nur die Störungen durch Jupiter berücksichtigt wurden. Letztere sind jedoch bei Weitem die stärksten. Ihre erste Berechnung hat mir ungefähr 400 Gleichungen ergeben. Bei einer zweiten, vollständigeren Rechnung, mit der ich bereits einen kleinen Anfang gemacht habe, werden ungefähr 1000 entwickelt, die über $0''1$ betragen. Die nächste Opposition wird nach der Rechnung von Herrn \textsc{Nicolai} eintreten

\Needspace{10\baselineskip}1813 August 18. $8^{\mathrm u}\;50^{\mathrm m}\;5^{\mathrm s}$ Mittlere Zeit in Göttingen

\begin{center}\fontsize{10}{12}\selectfont\setlength{\tabcolsep}{2pt}\renewcommand{\arraystretch}{1.15}\begin{tabular}{lcr}
Wahre Länge & $.\;.\;.\;.\;.\;.\;.\;.\;.\;.\;.\;.$ & $325^\circ\;24'\;18''$\\
Wahre geozentrische Breite & $.\;.\;.\;.\;.\;.\;.$ & $24\;37\;24\;\mathrm{N}..$\\
\end{tabular}\end{center}

\Needspace{17\baselineskip}\SourceLeaf{364}\par\indent Die Beobachtungen der \emph{Juno} waren dieses Jahr außerordentlich schwierig. Schon das bloße Auffinden in der von kleinen Sternen wimmelnden Milchstraße war sehr mühsam, weil der Planet so lichtschwach war. Er erreichte nur ungefähr die 11. und auch später noch nicht die 10. Größe. Als ich ihn aber erst einmal gefunden hatte, hatte ich das Glück, 4 gute Beobachtungen zu machen, die ich Ihnen hier mitteile.

\begin{center}\fontsize{10}{12}\selectfont\setlength{\tabcolsep}{4pt}\renewcommand{\arraystretch}{1.15}\begin{tabular}{l|c|c|c}
1812 & Mittlere Zeit & Scheinb. Rektasz. & Deklination\\\hline
Jan. 7 & $12^{\mathrm u}\;21^{\mathrm m}\;0^{\mathrm s}$ & $282^\circ\;21'\;24''5$ & $4^\circ\;41'\;44''6\;\mathrm S.$\\
13 & $12\;51\;12$ & $281\;16\;32.9$ & $4\;33\;39.6$\\
Juli 12 & $11\;35\;28$ & $275\;5\;10.3$ & $5\;6\;11.2$\\
13 & $11\;28\;12$ & $274\;52\;49.9$ & $5\;9\;20.0$\\
\end{tabular}\end{center}

\par\indent Aus diesen Beobachtungen hat Herr \textsc{Wachter} die Opposition nach einer eigenen Methode von mir wie folgt bestimmt. Sie erlaubt es, zeitlich weiter entfernte Beobachtungen einzubeziehen als das gewöhnliche Verfahren. Ihre Erklärung muss ich mir für eine andere Gelegenheit aufsparen:

\Needspace{10\baselineskip}1812 Juni 29. $19^{\mathrm u}\;51^{\mathrm m}\;43^{\mathrm s}$ Mittlere Zeit in Göttingen

\begin{center}\fontsize{10}{12}\selectfont\setlength{\tabcolsep}{2pt}\renewcommand{\arraystretch}{1.15}\begin{tabular}{lcr}
Wahre Länge & $.\;.\;.\;.\;.\;.\;.\;.\;.\;.\;.\;.$ & $278^\circ\;15'\;21''3$\\
Wahre geozentrische Breite & $.\;.\;.\;.\;.\;.\;.$ & $18\;36\;56.9$\\
\end{tabular}\end{center}

\Needspace{14\baselineskip}\par\indent Die Unterschiede zwischen meinen letzten Elementen und den Beobachtungen wurden auf die übliche Weise bestimmt. Nach Herrn \textsc{Wachter}s Rechnung sind es die folgenden:

\begin{center}\fontsize{10}{12}\selectfont\setlength{\tabcolsep}{7pt}\renewcommand{\arraystretch}{1.15}\begin{tabular}{l|c|c}
& Rektasz. & Dekl.\\\hline
Juni 7 & $+21''8$ & $-35''9$\\
13 & $+31.0$ & $-34.7$\\
Juli 12 & $+22.4$ & $-46.9$\\
13 & $+26.2$ & $-45\;0$\\
\end{tabular}\end{center}

\Needspace{16\baselineskip}\par\indent Derselbe geschickte Rechner hat außerdem die Elemente der 4 letzten Oppositionen nach meiner im I. Band der Neuen Comment. der hiesigen Königlichen Sozietät [S. 1 dieses Bandes] erklärten Methode angepasst und folgende Ergebnisse gefunden:

\begin{center}\fontsize{9}{11}\selectfont\setlength{\tabcolsep}{2pt}\renewcommand{\arraystretch}{1.15}\begin{tabular}{lcr}
Epoche, Göttinger Meridian & $.\;.\;.\;.\;.\;.\;.\;.$ & $177^\circ\;48'\;21''0$\\
Tägliche mittlere tropische Bewegung & $.\;.\;.\;.\;.$ & $813''25748$\\
Länge des Perihels 1811 & $.\;.\;.\;.\;.\;.\;.\;.\;.\;.\;.$ & $53^\circ\;15'\;10''1$\\
Länge des $\AscendingNode$ & $.\;.\;.\;.\;.\;.\;.\;.\;.\;.\;.\;.\;.\;.$ & $171\;9\;16.7$\\
Neigung der Bahn & $.\;.\;.\;.\;.\;.\;.\;.\;.\;.\;.\;.\;.$ & $13\;4\;17.2$\\
Logarithmus der großen Halbachse & $.\;.\;.\;.\;.\;.$ & $0.4265679$\\
\end{tabular}\end{center}

\par\indent Schließlich hat er auch die Ephemeride für Juno für das nächste Jahr berechnet. Ich habe das Vergnügen, Ihnen eine Kopie beizulegen. Im Übrigen scheint Juno in diesem Jahr nirgendwo sonst beobachtet worden zu sein.

\par\indent Herr Professor \textsc{Harding} hat Vesta ein paarmal am Kreismikrometer beobachtet. Die Ephemeride gab die Rektaszension um etwa $11'$ zu klein an. Ohne Zweifel wird ihre $\text{\AstroSymbols ☍}$ sehr gut durch Meridianbeobachtungen zu bestimmen sein. Vorher wollte ich dazu keine neuen Rechnungen vornehmen lassen. Es wird auch ausreichen, wenn eine neue Ephemeride etwa mit dem Ende des Jahres 1813 beginnt. Für das Jahrbuch 1815 wird also noch keine nötig sein.

\SourceLeaf{365}\par\indent Auch Ceres wurde zur Zeit der $\text{\AstroSymbols ☍}$ hier einigemal im und außerhalb des Meridians beobachtet. Dazu wurden jedoch noch keine Rechnungen angestellt.

\par\indent Professor \textsc{Harding} und ich haben einen Abend lang vergeblich nach dem in den Zeitungen angekündigten neuen Kometen gesucht.


